\section{A homogeneous symmetric kernel and the sign-only multiplicity case}
\label{sec:heisenberg-homogeneous-symmetric}

\paragraph{Author freeze and origin.}
D-TD41, reserved question Q2ew, is a derivation by
\texttt{/root/imprimitive} following the homogeneous-boundary task
of \texttt{/root}.  It retains the sign-only multiplicity-two
case explicitly.  It uses the actual normal-core obstruction
D-TD37 and the accepted finite-algebra and geometric-rank
arguments of D-TD27.  Complete author proof and source-interface
reading passes.  Independent and root whole-proof reviews of this new
written proof also pass.

Use the actual stable whole Heisenberg setup of D-TD05:
\[
 \begin{gathered}
 C\longrightarrow\mathbf P^1\text{ is a connected }J\text{-Galois cover},
 \qquad V=W\otimes R,\qquad e=\dim V,\\
 P=\operatorname{im}(J\to\operatorname{PGL}(W)),\qquad
 Q\triangleleft P,\quad Q\simeq\mathcal Q\simeq\mathbf F_p^{2m},
 \qquad P/Q=H\le\operatorname{Sp}(\mathcal Q),\qquad \dim W=p^m.
 \end{gathered}
 \tag{HSM.1}
\]
V is a faithful finite complex coefficient representation, the
projective Q-action is irreducible with nondegenerate commutator
form, and the original-field H-action on \(\mathcal Q\) is
faithful and irreducible.  R is the entire external tensor
partner, without a degree or support condition.  Suppose
\[
 \psi:H\twoheadrightarrow S_n,\quad n\ge6,\quad
 N=\ker\psi,\quad H_0=\psi^{-1}(A_n),\qquad
 \mathcal Q\big|_N\simeq S^{\oplus r},\quad r>1,
 \quad d=\dim_{\mathbf F_p}S.
 \tag{HSM.2}
\]
The restriction is semisimple: its nonzero N-socle is H-stable,
and original-field irreducibility makes it all of \(\mathcal Q\).
Let F be the finite commuting field of S, and write
\[
 F=\mathbf F_{p^f},\qquad d=af,\qquad
 \mathbf F_p[N]\big|_{\mathcal Q}
     =\operatorname{Mat}_a(F)\otimes1_r,
 \qquad \mathcal Q\simeq F^a\otimes_F F^r.
 \tag{HSM.3}
\]
These are image algebras in the original Fp representation.

\begin{theorem}[D-TD41]
There is an actual projective multiplicity homomorphism
\(\beta_0:A_n\to\operatorname{PGL}_r(F)\), defined below.
If it is nontrivial, it is faithful.  In this case the actual
moving-rank budget \(2e+6\), together with \(e>504\), forces
\[
 p^d<28900,\qquad d\le14,\qquad |N|<B:=28900^{14}.
 \tag{HSM.4}
\]
For \(n\ge B+2\), that case admits no actual sphere source
with the budget, without an additional \(e>504\) premise.

If \(\beta_0\) is trivial, H acts linearly over F and
\(r=2\).  Moreover p is odd and, for an odd natural element,
its multiplicity matrix B satisfies
\[
                   B^2=cI_2,\qquad c\in F^\times
                   \text{ is a nonsquare}.
 \tag{HSM.5}
\]
The exact full commuting field of H is then
\[
 F_H:=\operatorname{End}_H(\mathcal Q)=F[B],
 \qquad [F_H:F]=2,
 \qquad \dim_{F_H}\mathcal Q=a,
 \tag{HSM.6}
\]
and N is absolutely irreducible on this \(F_H\)-space.
This last possibility is described, not excluded.  In particular,
for \(n\ge B+2\), any source satisfying (HSM.1)--(HSM.2)
and the budget must lie in precisely this sign-only case.
\end{theorem}

\begin{proof}
\paragraph{Step 1: the actual field and projective multiplicity actions.}
The image-algebra statement (HSM.3) is the proved finite-algebra
argument of D-TD27.  H normalizes that algebra and its center F.
Conjugation by N acts trivially on its center.  Thus the field
action factors through \(S_n\) to the cyclic group
\(\operatorname{Gal}(F/\mathbf F_p)\).  It is either trivial
or the sign quotient.  H0 therefore centralizes F in both cases.

For \(h\in H_0\), Skolem--Noether and the full commuting
matrix algebra give
\[
 h=A(h)\otimes B(h),\qquad
 \beta_0:A_n\longrightarrow\operatorname{PGL}_r(F),
 \qquad \beta_0(hN)=[B(h)].
 \tag{HSM.7}
\]
The ambiguity in A and B consists of reciprocal F-scalars.
Consequently the projective map from H0 is a homomorphism;
it kills N and induces the displayed actual quotient map.
If \(\beta_0\) is nontrivial, simplicity of \(A_n\)
makes it faithful, and D-TD32 gives \(r>n/5\).

When the entire H action on F is trivial, the same construction
gives \(\beta:S_n\to\operatorname{PGL}_r(F)\) extending
\(\beta_0\).  If \(\beta_0\) is faithful then so is
\(\beta\): its kernel intersects \(A_n\) trivially, so
it centralizes \(A_n\), while the centralizer of \(A_n\)
in its natural \(S_n\) is trivial for \(n\ge5\).
This elementary natural centralizer fact is also used in D-TD37;
it makes no assertion about all automorphisms of \(A_6\).

When the H action on F is nontrivial, f is even.  Set
\(F_0=\mathbf F_{p^{f/2}}\), and let \(\sigma\) be
the nontrivial element of \(\operatorname{Gal}(F/F_0)\).
Every \(t\in H\setminus H_0\) is \(\sigma\)-semilinear.
Using entrywise \(\sigma\) in the chosen coordinates and
then the F-linear algebra-normalizer factorization gives
\[
 t=(A\otimes B)\sigma,
 \qquad A\in\operatorname{GL}_a(F),\quad
 B\in\operatorname{GL}_r(F).
 \tag{HSM.8}
\]
Here \(\sigma\) acts on both tensor factors, so the displayed
operator is semilinear on their tensor product over F.

\paragraph{Step 2: a fixed-space bound for every odd semilinear operator.}
Let T be any \(\sigma\)-semilinear invertible map on an
F-space of dimension s.  Its fixed vectors form an \(F_0\)-space.
Any \(F_0\)-independent set of fixed vectors is F-independent.
Indeed, a shortest nontrivial F-linear relation can be normalized
to have first coefficient one.  Apply T and subtract the original
relation.  Minimality forces every remaining coefficient to be
fixed by \(\sigma\), contradicting the \(F_0\)-independence.
It follows that
\[
 \dim_{\mathbf F_p}\ker(T-1)\le sf/2,
 \qquad \operatorname{rank}_{\mathbf F_p}(T-1)\ge sf/2.
 \tag{HSM.9}
\]
This argument does not assume \(T^2=1\), semisimplicity,
or an order prime to p.  Apply it to (HSM.8) with \(s=ar\).

\paragraph{Step 3: full-power support when the alternating action is nontrivial.}
Suppose \(\beta_0\ne1\).  If \(h\in H_0\) has
nonidentity natural image, B(h) is nonscalar.  Its geometric
projective rank over an algebraic closure is at least one.
D-TD27's tensor-rank inequality, multiplied by f, gives label
rank at least \(af=d\).  This remains true for nonsemisimple
operators.  For odd h, if the field action is trivial, the
faithfulness of \(\beta\) gives the same estimate.  If the
field action is nontrivial, (HSM.9) gives rank at least
\(rd/2\ge d\).  Thus
\[
 \psi(h)\ne1\quad\Longrightarrow\quad
 \operatorname{rank}_{\mathbf F_p}(1-h)\ge d,
 \qquad r>n/5.
 \tag{HSM.10}
\]
This applies to every actual power according to its own image.

Compose the actual maps \(J\to P\to H\to S_n\).
The connected symmetric quotient of the sphere cover has genus
at least two: \(n\ge6\), and the accepted genus-zero and
genus-one exclusions of D-TD02 apply to this non-solvable group
of order greater than sixty.  Its nonidentity inertia orders
therefore have orbifold sum at least \(85/42\).
The exact Heisenberg trace identity and (HSM.10) bound the
normalized absolute trace of every original element with
nonidentity natural image by \(c=p^{-d/2}\), with the whole
external partner retained.  For an original inertia order L
and its natural order o, exactly \(L/o\) powers have trivial
natural image.  Fourier inversion over the full original L
therefore gives every-target support at least
\((1-c)(1-1/o)\).  Consequently
\[
 2+6/e\ge(1-c)85/42,
 \qquad e(1-85c)\le252.
 \tag{HSM.11}
\]
For \(e>504\), this forces \(p^d<170^2=28900<2^{15}\).
N acts faithfully on S, because it acts faithfully on the sum
of copies \(S^{\oplus r}=\mathcal Q\).  Hence
\[
 |N|<p^{d^2}<28900^d\le B=28900^{14}.
 \tag{HSM.12}
\]
This proves (HSM.4).  No sphere-base assertion has been applied
to H0, to a restriction of the coefficient space, or to an
individual absolute constituent.

\paragraph{Step 4: the actual normal core for large n.}
For \(n\ge B+2\), (HSM.10) gives
\(e\ge p^{rd/2}\ge2^{r/2}>2^{n/10}>504\), since
\(B>10^6>100\).  Thus the presumed budget implies (HSM.12).
The induced conjugation map
\(A_n=H_0/N\to\operatorname{Out}(N)\) is trivial:
\(|\operatorname{Out}(N)|\le|N|!\le B!<n!/2=|A_n|\)
and \(A_n\) is simple.  Set
\[
 C_0=C_{H_0}(N),\qquad H_0=NC_0,\qquad
 C_0\cap N=Z(N),\qquad C_0/Z(N)\simeq A_n,
 \qquad E=C_0'.
 \tag{HSM.13}
\]
Both H0 and N are normal in H, so \(C_0=H_0\cap C_H(N)\)
and E are normal in all of H.  Since \(A_n\) is perfect,
\(C_0=E Z(N)\).  This central factorization gives
\(E=C_0'=E'\).  The image \(\psi(E)\) is \(A_n\),
and \(E\cap N=Z(E)\): the kernel lies in Z(N) and
centralizes E, while every central element of E has trivial
image in the centerless \(A_n\).  E is therefore an actual
normal quasisimple subgroup of the faithful label image.
All D-TD37 hypotheses now hold for the same original sphere
cover and entire partner, with \(n>10^6\).  Its conclusion
contradicts the budget \(2e+6\).

\paragraph{Step 5: a trivial alternating multiplicity map cannot be semilinear.}
Suppose \(\beta_0=1\) and the field action is nontrivial.
Fix an odd t and write it as (HSM.8).  On the multiplicity
space put \(T=B\sigma\).  Because \(t^2\in H_0\)
has scalar multiplicity action, uniqueness of tensor factors
implies \(T^2=cI_r\) for some \(c\in F^\times\).
The equality \(TT^2=T^2T\) gives \(\sigma(c)=c\),
so c lies in \(F_0^\times\).

The norm map \(F^\times\to F_0^\times\) is onto.
Indeed it is the power map with exponent \(|F_0|+1\),
its image lies in \(F_0^\times\), and its kernel has
at most \(|F_0|+1\) elements by the polynomial root bound.
Its image therefore has at least \(|F_0|-1\) elements and
is all of \(F_0^\times\).  Choose \(\lambda\in F^\times\) with
\(\lambda\sigma(\lambda)=c^{-1}\).  Then
\(T'=\lambda T\) is semilinear and \((T')^2=1\).
This reciprocal scalar change can also be absorbed in A, so it
does not change the original tensor operator t.

The fixed vectors of T' span the whole multiplicity space over F.
Here is an explicit descent argument, valid also in characteristic
two.  Choose \(u\in F\setminus F_0\).  For any w, both
\[
 s=w+T'w,\qquad z=uw+\sigma(u)T'w
 \quad\text{are fixed by }T',\qquad
 z-\sigma(u)s=(u-\sigma(u))w.
 \tag{HSM.14}
\]
The last scalar is nonzero.  There is therefore a nonzero fixed
vector v.  The F-line \(L=Fv\) is preserved by T', and
hence by T.  H0 has scalar multiplicity action, so
\(F^a\otimes_F L\) is preserved by H0 and by t.
Since H is generated by H0 and t and \(r>1\), this is a
nonzero proper original-Fp H-submodule, a contradiction.
Thus \(\beta_0=1\) forces the whole H action on F to be
trivial under the stated multiplicity hypothesis.

\paragraph{Step 6: the precise linear sign-only survivor.}
Now H is F-linear and \(\beta_0=1\).  Choose an odd t,
write \(t=A\otimes B\), and note that \(B^2=cI_r\)
because \(t^2\in H_0\).  Every H0 multiplicity matrix is
scalar.  Any proper nonzero B-invariant F-subspace L would make
\(F^a\otimes L\) a proper H-submodule, so B acts
irreducibly on \(F^r\).  An irreducible module for the
single-operator algebra \(F[B]\), whose defining polynomial
divides \(X^2-c\), has dimension at most two.  Therefore
\(r=2\).  The polynomial \(X^2-c\) must be irreducible.
In a finite field of characteristic two every element is a
square, and in odd characteristic this quadratic is irreducible
exactly when c is a nonsquare.  This proves (HSM.5).

Finally \(\operatorname{End}_N(\mathcal Q)=\operatorname{Mat}_2(F)\).
Every one of these matrices commutes with the scalar multiplicity
action of H0.  Commuting also with t means commuting with B.
The irreducible quadratic matrix has centralizer exactly F[B],
proving the first two assertions in (HSM.6).  This can also be
seen by identifying the multiplicity space with the one-dimensional
F[B]-space F[B]; its F-linear B-commuting maps are multiplication
by its elements.  Thus \(\mathcal Q\) has dimension a
over F[B].  The F[B]-span of the N image is
\[
          F[B][N]\big|_{\mathcal Q}
                      =\operatorname{Mat}_a(F[B]).
 \tag{HSM.15}
\]
It remains a full matrix algebra after algebraic closure, so N
is absolutely irreducible on this full-commuting-field model.
This proves the structural description without excluding it.
\end{proof}

\paragraph{Retained boundary.}
This theorem concerns homogeneous original-field N-restriction
of multiplicity greater than one.  It neither excludes the
sign-only case (HSM.5)--(HSM.6) nor the original multiplicity-one
case.  Reexpressing the sign-only case over its full commuting
field gives multiplicity one for N, and is not a source
contradiction.  No existence of a source, equality of complex
character fields, or polarization conclusion is inferred.
The index-two subgroup is used only for actual algebra and
normal-subgroup arguments; every geometric budget remains on
the full original symmetric quotient of the sphere cover.
